% Complete reviewed Unit 3; editorial master: chapters/ch03/answer-key.md
% Markdown SHA-256: 7a97bb18d56bb907e12abe76d82bf56b63d13211534928a58c1bc3d8709fc202
% One-use initial assembly followed by complete content and page review.
\newcommand{\runningtitle}{Unit 3. Answer key}
% Layout inherited from the reviewed previous course; set a running title per unit.
% This is a small hand-maintained TeX layout, not a conversion pipeline.
\RequirePackage{fix-cm}
\documentclass[11pt,twoside]{article}
\usepackage{fontspec}
\setmainfont{Cambria}[
  Renderer=HarfBuzz,
  BoldFont=cambriab.ttf,
  ItalicFont=cambriai.ttf,
  BoldItalicFont=cambriaz.ttf
]
\setmonofont{lmmono10-regular.otf}[BoldFont=lmmonolt10-bold.otf]
\usepackage[paperwidth=165mm,paperheight=241mm,inner=21mm,outer=19mm,
  top=22mm,bottom=20mm,headheight=14pt,headsep=7mm,footskip=10mm]{geometry}
\usepackage{xcolor,array,booktabs,longtable,tabularx}
\usepackage{enumitem,fancyhdr}
\usepackage[unicode,hidelinks]{hyperref}
\definecolor{ink}{HTML}{000000}
\definecolor{muted}{HTML}{4C5558}
\setlength{\parindent}{0pt}
\setlength{\parskip}{3.5pt}
\AtBeginDocument{\fontsize{11}{14.3}\selectfont}
\setlength{\tabcolsep}{4pt}
\renewcommand{\arraystretch}{1.10}
\setlist[itemize]{leftmargin=1.3em,label=\textbullet,itemsep=2pt,topsep=3pt,parsep=0pt}
\setlist[enumerate]{leftmargin=1.8em,itemsep=3pt,topsep=4pt,parsep=0pt}
\widowpenalty=10000
\clubpenalty=10000
\emergencystretch=1em
\raggedbottom
\pagestyle{fancy}
\fancyhf{}
\providecommand{\runningtitle}{Old English}
\fancyhead[LE,RO]{\small\itshape \runningtitle}
\fancyhead[LO,RE]{\small Old English}
\fancyfoot[LE,RO]{\small\thepage}
\renewcommand{\headrulewidth}{0.3pt}
\newcommand{\booktitle}[2]{
  {\fontsize{9}{11}\selectfont\color{muted}#1\par}
  \vspace{3pt}
  {\fontsize{20}{24}\selectfont\color{ink}#2\par}
  \vspace{6pt}
}
\newcommand{\source}[1]{\par{\footnotesize\color{muted}#1\par}}
\newcommand{\readingtext}[1]{
  \par\addvspace{5pt}
  \begingroup\leftskip=1.5em
    \fontsize{12}{16}\selectfont #1\par
  \endgroup\addvspace{5pt}
}
\makeatletter
\renewcommand{\section}{\@startsection{section}{1}{\z@}
  {-2.4ex \@plus -0.6ex \@minus -0.2ex}{0.8ex \@plus 0.2ex}
  {\normalfont\fontsize{12}{15}\selectfont\bfseries\color{ink}}}
\renewcommand{\subsection}{\@startsection{subsection}{2}{\z@}
  {-1.8ex \@plus -0.5ex \@minus -0.2ex}{0.5ex \@plus 0.2ex}
  {\normalfont\fontsize{11}{14}\selectfont\bfseries\color{ink}}}
\makeatother

\hyphenpenalty=10000
\exhyphenpenalty=10000
\emergencystretch=3em
\newcommand{\itempara}[2]{\par\addvspace{3pt}\noindent\hangindent=1.8em\hangafter=1\makebox[1.8em][l]{#1.}#2\par}
\newcommand{\reserve}[1]{\par\begingroup\dimen0=\pagegoal\advance\dimen0 by-\pagetotal\ifdim\dimen0<#1\newpage\fi\endgroup}
\begin{document}
\raggedright
\booktitle{UNIT 3}{Answer key}

\reserve{7\baselineskip}
% mdparagraph:P001:begin
\section*{Using the key}
% mdparagraph:P001:end

% mdparagraph:P002:begin
The models use the lesson’s selected forms. Documented spelling and inflectional alternatives remain valid Old English; distinguish them from recall of a specifically requested form. Word order may vary when the same cases, agreement, reference and meaning remain clear. Preserve vowel length and \textbf{ę}.
% mdparagraph:P002:end

% mdparagraph:P003:begin
Each row of a full paradigm labels every requested case, including identical spellings. Ordinary weak genitive plural \textbf{-ra} is the selected pattern, but Cook also admits \textbf{-ena} and \textbf{-a}. These alternatives are not errors merely because they differ from the model. Comparative stems need their own stated genitive-plural choices. An ordinary adjective after a possessive may be strong or weak; a possessive does not impose a universal weak ending. [1]
% mdparagraph:P003:end

% mdparagraph:P004:begin
Instrumental forms, \textbf{eall}, and plural \textbf{ān} have recognition/\allowbreak{}reference limits. Their display does not authorize independent instrumental syntax, quantifier constructions, or cardinal “one” with plural count nouns. The source-reading help remains available throughout its tasks. The answers and translations below are newly prepared models, not Cook’s own answer key.
% mdparagraph:P004:end

\reserve{7\baselineskip}
% mdparagraph:P005:begin
\section*{Adjectives}
% mdparagraph:P005:end

\par\begin{minipage}{\linewidth}\raggedright
% item: NEW-U03-A01
% key: KEY-NEW-U03-A01
% mdparagraph:P006:begin
\textbf{A1.} Complete strong \textbf{gōd}, “good”:
% mdparagraph:P006:end

% mdtable:T01:begin
\par\begingroup
\fontsize{10.5}{13.2}\selectfont
\setlength{\LTleft}{0pt}\setlength{\LTright}{\fill}
\noindent
\begin{tabular}{@{}>{\raggedright\arraybackslash}p{27.74mm}>{\raggedright\arraybackslash}p{21.25mm}>{\raggedright\arraybackslash}p{21.25mm}>{\raggedright\arraybackslash}p{21.25mm}>{\raggedright\arraybackslash}p{21.25mm}@{}}
\toprule
\textbf{Gender and number} & \textbf{Nominative} & \textbf{Genitive} & \textbf{Dative} & \textbf{Accusative} \\
\midrule
Masculine singular & gōd & gōdes & gōdum & gōdne \\
Neuter singular & gōd & gōdes & gōdum & gōd \\
Feminine singular & gōd & gōdre & gōdre & gōde \\
Masculine plural & gōde & gōdra & gōdum & gōde \\
Neuter plural & gōd & gōdra & gōdum & gōd \\
Feminine plural & gōda & gōdra & gōdum & gōda \\
\bottomrule
\end{tabular}
\par\endgroup

% mdtable:T01:end
% mdparagraph:P007:begin
The feminine nominative singular and neuter nominative/\allowbreak{}accusative plural have no \textbf{-u} after this long stem. Masculine plural \textbf{gōde} and feminine plural \textbf{gōda} remain distinct selected forms.
% mdparagraph:P007:end

\end{minipage}\par\addvspace{4pt}

\par\begin{minipage}{\linewidth}\raggedright
% item: NEW-U03-A02
% key: KEY-NEW-U03-A02
% mdparagraph:P008:begin
\textbf{A2.} Complete strong \textbf{glæd}, “glad”. Bold forms have stem \textbf{a} instead of \textbf{æ}:
% mdparagraph:P008:end

% mdtable:T02:begin
\par\begingroup
\fontsize{10.5}{13.2}\selectfont
\setlength{\LTleft}{0pt}\setlength{\LTright}{\fill}
\noindent
\begin{tabular}{@{}>{\raggedright\arraybackslash}p{27.74mm}>{\raggedright\arraybackslash}p{21.25mm}>{\raggedright\arraybackslash}p{21.25mm}>{\raggedright\arraybackslash}p{21.25mm}>{\raggedright\arraybackslash}p{21.25mm}@{}}
\toprule
\textbf{Gender and number} & \textbf{Nominative} & \textbf{Genitive} & \textbf{Dative} & \textbf{Accusative} \\
\midrule
Masculine singular & glæd & \textbf{glades} & \textbf{gladum} & glædne \\
Neuter singular & glæd & \textbf{glades} & \textbf{gladum} & glæd \\
Feminine singular & \textbf{gladu} & glædre & glædre & \textbf{glade} \\
Masculine plural & \textbf{glade} & glædra & \textbf{gladum} & \textbf{glade} \\
Neuter plural & \textbf{gladu} & glædra & \textbf{gladum} & \textbf{gladu} \\
Feminine plural & \textbf{glada} & glædra & \textbf{gladum} & \textbf{glada} \\
\bottomrule
\end{tabular}
\par\endgroup

% mdtable:T02:end
% mdparagraph:P009:begin
The selected vowel changes before the vowel endings shown here, including \textbf{glades} and \textbf{glade}. It remains \textbf{æ} before consonantal endings in \textbf{glædne, glædre, glædra}. Do not replace every \textbf{æ} mechanically.
% mdparagraph:P009:end

\end{minipage}\par\addvspace{4pt}

\par\begin{minipage}{\linewidth}\raggedright
% item: NEW-U03-A03
% key: KEY-NEW-U03-A03
% mdparagraph:P010:begin
\textbf{A3.} Complete weak \textbf{gōd}:
% mdparagraph:P010:end

% mdtable:T03:begin
\par\begingroup
\fontsize{10.5}{13.2}\selectfont
\setlength{\LTleft}{0pt}\setlength{\LTright}{\fill}
\noindent
\begin{tabular}{@{}>{\raggedright\arraybackslash}p{27.74mm}>{\raggedright\arraybackslash}p{21.25mm}>{\raggedright\arraybackslash}p{21.25mm}>{\raggedright\arraybackslash}p{21.25mm}>{\raggedright\arraybackslash}p{21.25mm}@{}}
\toprule
\textbf{Gender and number} & \textbf{Nominative} & \textbf{Genitive} & \textbf{Dative} & \textbf{Accusative} \\
\midrule
Masculine singular & gōda & gōdan & gōdan & gōdan \\
Neuter singular & gōde & gōdan & gōdan & gōde \\
Feminine singular & gōde & gōdan & gōdan & gōdan \\
Masculine plural & gōdan & gōdra & gōdum & gōdan \\
Neuter plural & gōdan & gōdra & gōdum & gōdan \\
Feminine plural & gōdan & gōdra & gōdum & gōdan \\
\bottomrule
\end{tabular}
\par\endgroup

% mdtable:T03:end
% mdparagraph:P011:begin
Selected adjective genitive plural \textbf{gōdra} contrasts with the ordinary weak-noun \textbf{-ena}, as in \textbf{steorrena}, “of stars”. Cook also allows weak-adjective \textbf{gōdena} and \textbf{gōda}: these are valid alternatives, not proof that the ordinary selected adjective ending must be \textbf{-ena}. [1]
% mdparagraph:P011:end

\end{minipage}\par\addvspace{4pt}

\par\begin{minipage}{\linewidth}\raggedright
% item: NEW-U03-A04
% key: KEY-NEW-U03-A04
% mdparagraph:P012:begin
\textbf{A4.} The requested cells are:
% mdparagraph:P012:end

% mdtable:T04:begin
\par\begingroup
\fontsize{10.5}{13.2}\selectfont
\setlength{\LTleft}{0pt}\setlength{\LTright}{\fill}
\noindent
\begin{tabular}{@{}>{\raggedright\arraybackslash}p{18.5mm}>{\raggedright\arraybackslash}p{23.75mm}>{\raggedright\arraybackslash}p{24.99mm}>{\raggedright\arraybackslash}p{22.75mm}>{\raggedright\arraybackslash}p{22.76mm}@{}}
\toprule
\textbf{Headword} & \textbf{Strong masculine genitive singular} & \textbf{Strong feminine nominative singular} & \textbf{Strong neuter nominative plural} & \textbf{Weak neuter nominative singular} \\
\midrule
micel & micles & miclu & miclu & micle \\
lȳtel & lȳtles & lȳtlu & lȳtlu & lȳtle \\
īdel & īdles & īdel & īdlu & īdle \\
ǣmtig & ǣmtiges & ǣmtig & ǣmtig & ǣmtige \\
\bottomrule
\end{tabular}
\par\endgroup

% mdtable:T04:end
% mdparagraph:P013:begin
The first three have selected \textbf{micl-, lȳtl-, īdl-} before the relevant vowel endings; \textbf{ǣmtig-} retains \textbf{i}. The endingless feminine \textbf{īdel/\allowbreak{}ǣmtig} and neuter plural \textbf{ǣmtig} are deliberate selections. Source-supported retained-vowel or endingless alternatives, such as feminine \textbf{lȳtelu/\allowbreak{}lȳtel}, are valid, although the task requests the selected forms.
% mdparagraph:P013:end

\end{minipage}\par\addvspace{4pt}

\par\begin{minipage}{\linewidth}\raggedright
% item: NEW-U03-A05
% key: KEY-NEW-U03-A05
% mdparagraph:P014:begin
\textbf{A5.} For \textbf{æppelbǣre}, “fruit-bearing”:
% mdparagraph:P014:end

% mdtable:T05:begin
\par\begingroup
\fontsize{10.5}{13.2}\selectfont
\setlength{\LTleft}{0pt}\setlength{\LTright}{\fill}
\noindent
\begin{tabular}{@{}>{\raggedright\arraybackslash}p{82.96mm}>{\raggedright\arraybackslash}p{38.23mm}@{}}
\toprule
\textbf{Cell} & \textbf{Form} \\
\midrule
Strong masculine nominative singular & æppelbǣre \\
Strong feminine nominative singular & æppelbǣru \\
Strong masculine accusative singular & æppelbǣrne \\
Strong neuter nominative plural & æppelbǣru \\
Strong genitive plural, all genders & æppelbǣrra \\
\bottomrule
\end{tabular}
\par\endgroup

% mdtable:T05:end
% mdparagraph:P015:begin
The headword belongs to the \textbf{-e} pattern. Its \textbf{-e} is present where no other ending replaces it; it is not evidence that the masculine nominative is weak. Adding the consonantal endings gives \textbf{-rne/\allowbreak{}-rra}, not an extra retained \textbf{e}.
% mdparagraph:P015:end

% mdtable:T06:begin
\par\begingroup
\fontsize{10}{12.8}\selectfont
\setlength{\LTleft}{0pt}\setlength{\LTright}{\fill}
\noindent
\begin{tabular}{@{}>{\raggedright\arraybackslash}p{18.5mm}>{\raggedright\arraybackslash}p{19.39mm}>{\raggedright\arraybackslash}p{17.69mm}>{\raggedright\arraybackslash}p{16.63mm}>{\raggedright\arraybackslash}p{19.39mm}>{\raggedright\arraybackslash}p{18.33mm}@{}}
\toprule
\textbf{Headword} & \textbf{Strong masculine nominative} & \textbf{Strong masculine genitive} & \textbf{Strong masculine dative} & \textbf{Strong masculine accusative} & \textbf{Weak masculine nominative} \\
\midrule
gearu & gearu & gearwes & gearwum & gearone & gearwa \\
hēah & hēah & hēas & hēaum & hēanne & hēa \\
\bottomrule
\end{tabular}
\par\endgroup

% mdtable:T06:end
% mdparagraph:P016:begin
\textbf{Gearu} has \textbf{gearw-} before vowel endings and selected \textbf{gearo-} before the consonantal accusative ending. \textbf{Hēah} loses \textbf{h} and contracts in its selected oblique forms. Dative \textbf{hēam} and accusative \textbf{hēane} are documented alternatives; they do not change the case analysis.
% mdparagraph:P016:end

\end{minipage}\par\addvspace{4pt}

\par\begin{minipage}{\linewidth}\raggedright
% item: NEW-U03-A06
% key: KEY-NEW-U03-A06
% mdparagraph:P017:begin
\textbf{A6.} \textbf{Ðing}, neuter, “thing”:
% mdparagraph:P017:end

% mdtable:T07:begin
\par\begingroup
\fontsize{10.5}{13.2}\selectfont
\setlength{\LTleft}{0pt}\setlength{\LTright}{\fill}
\noindent
\begin{tabular}{@{}>{\raggedright\arraybackslash}p{39.46mm}>{\raggedright\arraybackslash}p{39.46mm}>{\raggedright\arraybackslash}p{39.46mm}@{}}
\toprule
\textbf{Case} & \textbf{Singular} & \textbf{Plural} \\
\midrule
Nominative & ðing & ðing \\
Genitive & ðinges & ðinga \\
Dative & ðinge & ðingum \\
Accusative & ðing & ðing \\
\bottomrule
\end{tabular}
\par\endgroup

% mdtable:T07:end
% mdparagraph:P018:begin
\textbf{Ðā ðing} means “the things”. The plural nominative/\allowbreak{}accusative has no added ending because this is a heavy-stem neuter. Identical singular and plural noun spellings do not make the phrase singular.
% mdparagraph:P018:end

\end{minipage}\par\addvspace{4pt}

\par\begin{minipage}{\linewidth}\raggedright
% item: NEW-U03-A07
% key: KEY-NEW-U03-A07
% mdparagraph:P019:begin
\textbf{A7.} \textbf{Se gōda fugol; sēo gōde wyrt; ðæt gōde trēo; ðā gōdan fuglas.} “The good bird; the good herb; the good tree; the good birds.” Each adjective becomes weak after its article. The first three remain singular, the fourth masculine plural; all remain nominative.
% mdparagraph:P019:end

\end{minipage}\par\addvspace{4pt}

\par\begin{minipage}{\linewidth}\raggedright
% item: NEW-U03-A08
% key: KEY-NEW-U03-A08
% mdparagraph:P020:begin
\textbf{A8.} “The good birds are in the dwelling. The birds are good.” \textbf{Gōdan} is weak nominative masculine plural inside \textbf{ðā gōdan fuglas}. Predicate \textbf{gōde} is strong nominative masculine plural, agreeing with \textbf{fuglas}. The subject’s article does not make a separate predicate adjective weak. \textbf{On ðǣm hofe} is a dative location.
% mdparagraph:P020:end

\end{minipage}\par\addvspace{4pt}

\par\begin{minipage}{\linewidth}\raggedright
% item: NEW-U03-A09
% key: KEY-NEW-U03-A09
% mdparagraph:P021:begin
\textbf{A9.} \textbf{Ðā fuglas libbað on ðæs gōdan mannes hofe.} The birds are the nominative plural subject. \textbf{Hofe} is dative neuter singular after locative \textbf{on}; \textbf{ðæs gōdan mannes} is an embedded genitive masculine singular phrase, “of the good man”. \textbf{Gōdan} agrees with \textbf{mannes}, not \textbf{hofe}. Placing the genitive after the head can express the same relation, but does not satisfy the requested position.
% mdparagraph:P021:end

\end{minipage}\par\addvspace{4pt}

\par\begin{minipage}{\linewidth}\raggedright
% item: NEW-U03-A10
% key: KEY-NEW-U03-A10
% mdparagraph:P022:begin
\textbf{A10.} \textbf{Gōde} can be strong masculine nominative plural, masculine instrumental singular, or neuter instrumental singular in the stated table cells. The corresponding \textbf{glæd} forms are \textbf{glade} in all three cells, with \textbf{a}. Identical spellings do not collapse the analyses. Recognition of the instrumental cells is sufficient; no independent instrumental construction is required.
% mdparagraph:P022:end

\end{minipage}\par\addvspace{4pt}

\reserve{7\baselineskip}
% mdparagraph:P023:begin
\section*{Possessives}
% mdparagraph:P023:end

\par\begin{minipage}{\linewidth}\raggedright
% item: NEW-U03-P01
% key: KEY-NEW-U03-P01
% mdparagraph:P024:begin
\textbf{P1.} Complete agreeing \textbf{mīn}, “my”:
% mdparagraph:P024:end

% mdtable:T08:begin
\par\begingroup
\fontsize{10.5}{13.2}\selectfont
\setlength{\LTleft}{0pt}\setlength{\LTright}{\fill}
\noindent
\begin{tabular}{@{}>{\raggedright\arraybackslash}p{27.74mm}>{\raggedright\arraybackslash}p{21.25mm}>{\raggedright\arraybackslash}p{21.25mm}>{\raggedright\arraybackslash}p{21.25mm}>{\raggedright\arraybackslash}p{21.25mm}@{}}
\toprule
\textbf{Gender and number} & \textbf{Nominative} & \textbf{Genitive} & \textbf{Dative} & \textbf{Accusative} \\
\midrule
Masculine singular & mīn & mīnes & mīnum & mīnne \\
Neuter singular & mīn & mīnes & mīnum & mīn \\
Feminine singular & mīn & mīnre & mīnre & mīne \\
Masculine plural & mīne & mīnra & mīnum & mīne \\
Neuter plural & mīn & mīnra & mīnum & mīn \\
Feminine plural & mīna & mīnra & mīnum & mīna \\
\bottomrule
\end{tabular}
\par\endgroup

% mdtable:T08:end
% mdparagraph:P025:begin
The owner is first-person singular throughout. The \textbf{possessed noun’s gender, number and case} determine the ending. Thus a single owner can possess a plural noun; the possessive ending does not count the owners.
% mdparagraph:P025:end

\end{minipage}\par\addvspace{4pt}

\par\begin{minipage}{\linewidth}\raggedright
% item: NEW-U03-P02
% key: KEY-NEW-U03-P02
% mdparagraph:P026:begin
\textbf{P2.} \textbf{Ic hīere ðīnne fugol. Ðū hīerst mīne fuglas.} \textbf{Ðīnne} is masculine accusative singular with \textbf{fugol}; \textbf{mīne} is masculine accusative plural with \textbf{fuglas}. The singular addressee selects \textbf{ðīn-} and \textbf{ðū}, while the two subjects select \textbf{hīere/\allowbreak{}hīerst}. Preserve the doubled \textbf{nn} in \textbf{ðīnne}.
% mdparagraph:P026:end

\end{minipage}\par\addvspace{4pt}

\par\begin{minipage}{\linewidth}\raggedright
% item: NEW-U03-P03
% key: KEY-NEW-U03-P03
% mdparagraph:P027:begin
\textbf{P3.} \textbf{Wit habbað hof. Uncer hof is lȳtel. Git habbað trēo. Incer trēo is micel.} \textbf{Wit/\allowbreak{}git} refer to two people and take plural \textbf{habbað}. \textbf{Uncer/\allowbreak{}incer} preserve the two-person possessors, but each agrees with a singular neuter noun. Consequently \textbf{hof/\allowbreak{}trēo} take singular \textbf{is} and strong neuter singular predicates \textbf{lȳtel/\allowbreak{}micel}.
% mdparagraph:P027:end

\end{minipage}\par\addvspace{4pt}

\par\begin{minipage}{\linewidth}\raggedright
% item: NEW-U03-P04
% key: KEY-NEW-U03-P04
% mdparagraph:P028:begin
\textbf{P4.} \textbf{Ūre fuglas libbað on ēowrum hofe.} \textbf{Ūre} identifies several speakers as owners and agrees with masculine nominative plural \textbf{fuglas}. \textbf{Ēowrum} identifies several addressees as owners and agrees with neuter dative singular \textbf{hofe}. The locative construction, not the number of owners, determines the dative.
% mdparagraph:P028:end

\end{minipage}\par\addvspace{4pt}

\par\begin{minipage}{\linewidth}\raggedright
% item: NEW-U03-P05
% key: KEY-NEW-U03-P05
% mdparagraph:P029:begin
\textbf{P5.} \textbf{His fuglas libbað on mīnum hofe. Hiera hof is māre ðonne ūre hof.} \textbf{His/\allowbreak{}hiera} are invariant third-person possessors, here a male singular owner and plural owners. \textbf{Mīnum} agrees with dative neuter singular \textbf{hofe}; \textbf{ūre} with nominative neuter singular \textbf{hof}. Both dwellings in the comparison are nominative subjects, with the second \textbf{is} unexpressed. Adding that repeated \textbf{is} is acceptable.
% mdparagraph:P029:end

\end{minipage}\par\addvspace{4pt}

\par\begin{minipage}{\linewidth}\raggedright
% item: NEW-U03-P06
% key: KEY-NEW-U03-P06
% mdparagraph:P030:begin
\textbf{P6.} The given contexts fix these analyses:
% mdparagraph:P030:end

% mdtable:T09:begin
\par\begingroup
\fontsize{10.5}{13.2}\selectfont
\setlength{\LTleft}{0pt}\setlength{\LTright}{\fill}
\noindent
\begin{tabular}{@{}>{\raggedright\arraybackslash}p{25.99mm}>{\raggedright\arraybackslash}p{35.99mm}>{\raggedright\arraybackslash}p{56.39mm}@{}}
\toprule
\textbf{Phrase} & \textbf{Owner} & \textbf{Possessed noun} \\
\midrule
mīne fuglas & first person singular & fuglas: masculine nominative plural \\
uncer fugol & first person dual & fugol: masculine nominative singular \\
ēowrum hofe & second person plural & hofe: neuter dative singular \\
\bottomrule
\end{tabular}
\par\endgroup

% mdtable:T09:end
% mdparagraph:P031:begin
The stem identifies the owner’s person and number. The inflected form agrees with the possessed noun; its ending may coincide with other gender/\allowbreak{}number/\allowbreak{}case cells. Here the subject or locative context and the noun establish the precise analysis. \textbf{Uncer} has no added ending in this nominative cell; it does not make the bird dual.
% mdparagraph:P031:end

\end{minipage}\par\addvspace{4pt}

\reserve{7\baselineskip}
% mdparagraph:P032:begin
\section*{Comparison}
% mdparagraph:P032:end

\par\begin{minipage}{\linewidth}\raggedright
% item: NEW-U03-C01
% key: KEY-NEW-U03-C01
% mdparagraph:P033:begin
\textbf{C1.} Selected masculine nominative singular forms:
% mdparagraph:P033:end

% mdtable:T10:begin
\par\begingroup
\fontsize{10.5}{13.2}\selectfont
\setlength{\LTleft}{0pt}\setlength{\LTright}{\fill}
\noindent
\begin{tabular}{@{}>{\raggedright\arraybackslash}p{34.56mm}>{\raggedright\arraybackslash}p{38.02mm}>{\raggedright\arraybackslash}p{45.8mm}@{}}
\toprule
\textbf{Positive} & \textbf{Comparative} & \textbf{Weak superlative} \\
\midrule
lēof & lēofra & lēofosta \\
eald & ieldra & ieldesta \\
gōd & bętera & bętsta \\
micel & māra & mǣsta \\
lȳtel & lǣssa & lǣsta \\
yfel & wiersa & wierresta \\
\bottomrule
\end{tabular}
\par\endgroup

% mdtable:T10:end
% mdparagraph:P034:begin
\textbf{Lēof} illustrates regular comparison; \textbf{eald} has the selected umlaut forms. The remaining four series require their lexical stems. A strong superlative citation such as \textbf{bętst} is a real form, but it does not meet the request for a weak masculine nominative singular.
% mdparagraph:P034:end

\end{minipage}\par\addvspace{4pt}

\par\begin{minipage}{\linewidth}\raggedright
% item: NEW-U03-C02
% key: KEY-NEW-U03-C02
% mdparagraph:P035:begin
\textbf{C2.} Complete weak comparative \textbf{māra}, “greater/\allowbreak{}larger”:
% mdparagraph:P035:end

% mdtable:T11:begin
\par\begingroup
\fontsize{10.5}{13.2}\selectfont
\setlength{\LTleft}{0pt}\setlength{\LTright}{\fill}
\noindent
\begin{tabular}{@{}>{\raggedright\arraybackslash}p{27.74mm}>{\raggedright\arraybackslash}p{21.25mm}>{\raggedright\arraybackslash}p{21.25mm}>{\raggedright\arraybackslash}p{21.25mm}>{\raggedright\arraybackslash}p{21.25mm}@{}}
\toprule
\textbf{Gender and number} & \textbf{Nominative} & \textbf{Genitive} & \textbf{Dative} & \textbf{Accusative} \\
\midrule
Masculine singular & māra & māran & māran & māran \\
Neuter singular & māre & māran & māran & māre \\
Feminine singular & māre & māran & māran & māran \\
Masculine plural & māran & mārra & mārum & māran \\
Neuter plural & māran & mārra & mārum & māran \\
Feminine plural & māran & mārra & mārum & māran \\
\bottomrule
\end{tabular}
\par\endgroup

% mdtable:T11:end
% mdparagraph:P036:begin
The neuter nominative/\allowbreak{}accusative singular is \textbf{māre}, and the selected genitive plural is \textbf{mārra}, with \textbf{rr}. In the six comparison series the selected genitive plurals are \textbf{lēofra, ieldra, bęterra, mārra, lǣssena, wiersena}. The first two simplify the repeated consonant after another consonant; the latter two use the licensed \textbf{-ena} alternative. Do not form every comparative genitive plural by mechanically attaching written \textbf{-ra}. The selected completions apply checked rules; a defensible source variant remains valid language. [1]
% mdparagraph:P036:end

\end{minipage}\par\addvspace{4pt}

\par\begin{minipage}{\linewidth}\raggedright
% item: NEW-U03-C03
% key: KEY-NEW-U03-C03
% mdparagraph:P037:begin
\textbf{C3.} For \textbf{mǣst}, “greatest”:
% mdparagraph:P037:end

% mdtable:T12:begin
\par\begingroup
\fontsize{10}{12.8}\selectfont
\setlength{\LTleft}{0pt}\setlength{\LTright}{\fill}
\noindent
\begin{tabular}{@{}>{\raggedright\arraybackslash}p{19.2mm}>{\raggedright\arraybackslash}p{19.60mm}>{\raggedright\arraybackslash}p{18.36mm}>{\raggedright\arraybackslash}p{18.78mm}>{\raggedright\arraybackslash}p{17.83mm}>{\raggedright\arraybackslash}p{16.17mm}@{}}
\toprule
\textbf{Declension} & \textbf{Masculine nominative singular} & \textbf{Neuter nominative singular} & \textbf{Feminine nominative singular} & \textbf{Genitive plural, all genders} & \textbf{Dative plural, all genders} \\
\midrule
Strong & mǣst & mǣst & mǣst & mǣstra & mǣstum \\
Weak & mǣsta & mǣste & mǣste & mǣstra & mǣstum \\
\bottomrule
\end{tabular}
\par\endgroup

% mdtable:T12:end
% mdparagraph:P038:begin
The degree remains superlative in both declensions. The coinciding plural obliques do not erase the strong/\allowbreak{}weak distinction in the singular. The table gives the selected genitive plural; the ordinary weak alternatives described above remain possible.
% mdparagraph:P038:end

\end{minipage}\par\addvspace{4pt}

\par\begin{minipage}{\linewidth}\raggedright
% item: NEW-U03-C04
% key: KEY-NEW-U03-C04
% mdparagraph:P039:begin
\textbf{C4.} \textbf{Ðæt micle trēo is māre ðonne ðæt lȳtle trēo.} Both trees are nominative neuter singular: they are the contrasted subjects, with a repeated \textbf{is} understood after the second tree. \textbf{Micle/\allowbreak{}lȳtle} are weak attributive adjectives after articles. \textbf{Māre} is weak because it is comparative, and \textbf{-e} is its neuter nominative singular ending. A repeated final \textbf{is} is acceptable. [2]
% mdparagraph:P039:end

\end{minipage}\par\addvspace{4pt}

\par\begin{minipage}{\linewidth}\raggedright
% item: NEW-U03-C05
% key: KEY-NEW-U03-C05
% mdparagraph:P040:begin
\textbf{C5.} \textbf{Se fugol is lēofra ðonne se fisc. Ðæt ealde hof is bętere ðonne ðæt micle hof.} The first comparative is masculine nominative singular, the second neuter nominative singular. Both are weak despite being predicates. The standards \textbf{se fisc} and \textbf{ðæt micle hof} are nominative in these parallel comparisons. The second sentence requires adjectival \textbf{bętere}, not adverbial \textbf{bęt}.
% mdparagraph:P040:end

\end{minipage}\par\addvspace{4pt}

\par\begin{minipage}{\linewidth}\raggedright
% item: NEW-U03-C06
% key: KEY-NEW-U03-C06
% mdparagraph:P041:begin
\textbf{C6.} \textbf{Ðæt ieldeste trēo; se lēofosta fugol; ðæt mǣste weorc; se bętsta mann; se lǣsta steorra.} All are nominative singular with weak superlatives. \textbf{Trēo/\allowbreak{}weorc} are neuter and take \textbf{-e}; \textbf{fugol/\allowbreak{}mann/\allowbreak{}steorra} are masculine and take \textbf{-a}. Preserve \textbf{ield-}, with no added macron, and the irregular superlative stems.
% mdparagraph:P041:end

\end{minipage}\par\addvspace{4pt}

\par\begin{minipage}{\linewidth}\raggedright
% item: NEW-U03-C07
% key: KEY-NEW-U03-C07
% mdparagraph:P042:begin
\textbf{C7.} \textbf{Se fugol is lǣssa. Ðā fuglas sindon lǣssan.} “The bird is smaller. The birds are smaller.” Comparative \textbf{lǣssa} is weak masculine nominative singular; the plural is \textbf{lǣssan}. Bare \textbf{lǣss} does not supply the required agreeing comparative in the selected system. The subject’s number also changes \textbf{is} to \textbf{sindon}.
% mdparagraph:P042:end

\end{minipage}\par\addvspace{4pt}

\reserve{7\baselineskip}
% mdparagraph:P043:begin
\section*{Adverbs}
% mdparagraph:P043:end

\par\begin{minipage}{\linewidth}\raggedright
% item: NEW-U03-V01
% key: KEY-NEW-U03-V01
% mdparagraph:P044:begin
\textbf{V1.} \textbf{Wīde, wīdor, wīdost}: “widely, more widely, most widely”. From supplied \textbf{sōð}, the selected adverb is \textbf{sōðlice}, “truly”. None changes for gender, case or number. \textbf{Wīd/\allowbreak{}sōð} are supplied derivational citations here; the task does not require their independent adjective paradigms. Do not add a macron to the \textbf{i} of \textbf{-lice}.
% mdparagraph:P044:end

\end{minipage}\par\addvspace{4pt}

\par\begin{minipage}{\linewidth}\raggedright
% item: NEW-U03-V02
% key: KEY-NEW-U03-V02
% mdparagraph:P045:begin
\textbf{V2.} \textbf{Ðā męnn geworhton ðæt gōde hof wel.} \textbf{Gōde} is a weak accusative neuter singular adjective modifying the dwelling \textbf{hof}. \textbf{Wel} is an indeclinable adverb modifying the building action \textbf{geworhton}. The men control the plural verb; the adverb does not agree with either noun.
% mdparagraph:P045:end

\end{minipage}\par\addvspace{4pt}

\par\begin{minipage}{\linewidth}\raggedright
% item: NEW-U03-V03
% key: KEY-NEW-U03-V03
% mdparagraph:P046:begin
\textbf{V3.} \textbf{Hē geworhte ðæt hof bęt. Hē geworhte ðæt hof bętst.} “He built the dwelling better/\allowbreak{}best.” \textbf{Bęt/\allowbreak{}bętst} are the comparative and superlative adverbs of \textbf{wel}. They qualify how he built, so they have no agreement with accusative neuter \textbf{hof}. Adjectival \textbf{bętere} would not fulfill the requested adverbial target.
% mdparagraph:P046:end

\end{minipage}\par\addvspace{4pt}

\par\begin{minipage}{\linewidth}\raggedright
% item: NEW-U03-V04
% key: KEY-NEW-U03-V04
% mdparagraph:P047:begin
\textbf{V4.} “Truly, the work is very good”, compared with “The work is good”. \textbf{Sōðlice} presents the statement as true/\allowbreak{}indeed so; \textbf{swīðe} increases the degree of \textbf{gōd}. The adverbs do not inflect for the noun; predicate \textbf{gōd} does agree with neuter nominative singular \textbf{weorc}. This “truly” sense is distinct from the locally supplied narrative connective \textbf{sōðlice} in the Creation passage.
% mdparagraph:P047:end

\end{minipage}\par\addvspace{4pt}

\reserve{7\baselineskip}
% mdparagraph:P048:begin
\section*{Cardinals}
% mdparagraph:P048:end

\par\begin{minipage}{\linewidth}\raggedright
% item: NEW-U03-N01
% key: KEY-NEW-U03-N01
% mdparagraph:P049:begin
\textbf{N1.} Selected singular cardinal \textbf{ān}, “one”:
% mdparagraph:P049:end

% mdtable:T13:begin
\par\begingroup
\fontsize{10.5}{13.2}\selectfont
\setlength{\LTleft}{0pt}\setlength{\LTright}{\fill}
\noindent
\begin{tabular}{@{}>{\raggedright\arraybackslash}p{28.89mm}>{\raggedright\arraybackslash}p{28.89mm}>{\raggedright\arraybackslash}p{28.89mm}>{\raggedright\arraybackslash}p{28.89mm}@{}}
\toprule
\textbf{Case} & \textbf{Masculine} & \textbf{Neuter} & \textbf{Feminine} \\
\midrule
Nominative & ān & ān & ān \\
Genitive & ānes & ānes & ānre \\
Dative & ānum & ānum & ānre \\
Accusative & ānne & ān & āne \\
\bottomrule
\end{tabular}
\par\endgroup

% mdtable:T13:end
% mdparagraph:P050:begin
The number is one, so the accompanying count noun is singular. Preserve masculine accusative \textbf{ānne}. The reference plural forms do not license translating cardinal “one” with a plural count noun.
% mdparagraph:P050:end

\end{minipage}\par\addvspace{4pt}

\par\begin{minipage}{\linewidth}\raggedright
% item: NEW-U03-N02
% key: KEY-NEW-U03-N02
% mdparagraph:P051:begin
\textbf{N2.} The complete selected plural systems are:
% mdparagraph:P051:end

% mdparagraph:P052:begin
\textbf{Twēgen}, “two”:
% mdparagraph:P052:end

% mdtable:T14:begin
\par\begingroup
\fontsize{10.5}{13.2}\selectfont
\setlength{\LTleft}{0pt}\setlength{\LTright}{\fill}
\noindent
\begin{tabular}{@{}>{\raggedright\arraybackslash}p{28.89mm}>{\raggedright\arraybackslash}p{28.89mm}>{\raggedright\arraybackslash}p{28.89mm}>{\raggedright\arraybackslash}p{28.89mm}@{}}
\toprule
\textbf{Case} & \textbf{Masculine} & \textbf{Neuter} & \textbf{Feminine} \\
\midrule
Nominative & twēgen & twā & twā \\
Genitive & twēgra & twēgra & twēgra \\
Dative & twǣm & twǣm & twǣm \\
Accusative & twēgen & twā & twā \\
\bottomrule
\end{tabular}
\par\endgroup

% mdtable:T14:end
% mdparagraph:P053:begin
\textbf{Ðrīe}, “three”:
% mdparagraph:P053:end

% mdtable:T15:begin
\par\begingroup
\fontsize{10.5}{13.2}\selectfont
\setlength{\LTleft}{0pt}\setlength{\LTright}{\fill}
\noindent
\begin{tabular}{@{}>{\raggedright\arraybackslash}p{28.89mm}>{\raggedright\arraybackslash}p{28.89mm}>{\raggedright\arraybackslash}p{28.89mm}>{\raggedright\arraybackslash}p{28.89mm}@{}}
\toprule
\textbf{Case} & \textbf{Masculine} & \textbf{Neuter} & \textbf{Feminine} \\
\midrule
Nominative & ðrīe & ðrēo & ðrēo \\
Genitive & ðrēora & ðrēora & ðrēora \\
Dative & ðrīm & ðrīm & ðrīm \\
Accusative & ðrīe & ðrēo & ðrēo \\
\bottomrule
\end{tabular}
\par\endgroup

% mdtable:T15:end
% mdparagraph:P054:begin
Both modify \textbf{plural} count nouns. Two does not require a special dual noun declension. Cook’s \textbf{twēga}, \textbf{twām}, and neuter \textbf{tū} are valid alternatives to selected \textbf{twēgra}, \textbf{twǣm}, and \textbf{twā} in their respective cells.
% mdparagraph:P054:end

\end{minipage}\par\addvspace{4pt}

\par\begin{minipage}{\linewidth}\raggedright
% item: NEW-U03-N03
% key: KEY-NEW-U03-N03
% mdparagraph:P055:begin
\textbf{N3.} \textbf{Twēgen męnn; twā wyrta; twā hofu; ðrīe fuglas; ðrēo niht; ðrēo trēowu.} All six phrases are nominative plural. The two numerals distinguish masculine from neuter/\allowbreak{}feminine forms. Plural \textbf{niht} has no added ending here; it still means three nights. Neuter \textbf{tū hofu} is a documented alternative to \textbf{twā hofu}.
% mdparagraph:P055:end

\end{minipage}\par\addvspace{4pt}

\par\begin{minipage}{\linewidth}\raggedright
% item: NEW-U03-N04
% key: KEY-NEW-U03-N04
% mdparagraph:P056:begin
\textbf{N4.} \textbf{Wē sealdon twǣm mannum sǣd. Hīe sealdon ðrīm fuglum wæter.} The numeral and recipient noun are dative plural in each sentence; \textbf{sǣd/\allowbreak{}wæter} are the accusative things given. \textbf{Twām} is a valid dative alternative. Nominative \textbf{twēgen męnn} or \textbf{ðrīe fuglas} would not express the recipients in this construction.
% mdparagraph:P056:end

\end{minipage}\par\addvspace{4pt}

\par\begin{minipage}{\linewidth}\raggedright
% item: NEW-U03-N05
% key: KEY-NEW-U03-N05
% mdparagraph:P057:begin
\textbf{N5.} The selected nominative phrases are:
% mdparagraph:P057:end

% mdtable:T16:begin
\par\begingroup
\fontsize{10.5}{13.2}\selectfont
\setlength{\LTleft}{0pt}\setlength{\LTright}{\fill}
\noindent
\begin{tabular}{@{}>{\raggedright\arraybackslash}p{59.85mm}>{\raggedright\arraybackslash}p{61.34mm}@{}}
\toprule
\textbf{Meaning} & \textbf{Old English} \\
\midrule
four fish & fēower fiscas \\
five trees & fīf trēowu \\
six stars & siex steorran \\
seven days & seofon dagas \\
eight birds & eahta fuglas \\
nine men & nigon męnn \\
ten herbs & tīen wyrta \\
eleven years & endlefan gēar \\
twelve signs & twęlf tācnu \\
\bottomrule
\end{tabular}
\par\endgroup

% mdtable:T16:end
% mdparagraph:P058:begin
The nouns are plural count nouns; the cardinals four through twelve are uninflected in this taught position before them. \textbf{Gēar} has a zero plural ending, and \textbf{tācnu} shows noun-stem syncope. Do not replace these plural noun forms with singulars merely because the cardinals are invariant. Retain selected \textbf{endlefan} and hooked \textbf{ę} in \textbf{twęlf}.
% mdparagraph:P058:end

\end{minipage}\par\addvspace{4pt}

\reserve{7\baselineskip}
% mdparagraph:P059:begin
\section*{Ordinals}
% mdparagraph:P059:end

\par\begin{minipage}{\linewidth}\raggedright
% item: NEW-U03-O01
% key: KEY-NEW-U03-O01
% mdparagraph:P060:begin
\textbf{O1.} The singular phrases are:
% mdparagraph:P060:end

% mdtable:T17:begin
\par\begingroup
\fontsize{10.5}{13.2}\selectfont
\setlength{\LTleft}{0pt}\setlength{\LTright}{\fill}
\noindent
\begin{tabular}{@{}>{\raggedright\arraybackslash}p{38.47mm}>{\raggedright\arraybackslash}p{39.52mm}>{\raggedright\arraybackslash}p{40.39mm}@{}}
\toprule
\textbf{Meaning} & \textbf{Nominative} & \textbf{Genitive} \\
\midrule
the second day & se ōðer dæg & ðæs ōðres dæges \\
the second night & sēo ōðru niht & ðǣre ōðerre niht \\
the second year & ðæt ōðer gēar & ðæs ōðres gēares \\
\bottomrule
\end{tabular}
\par\endgroup

% mdtable:T17:end
% mdparagraph:P061:begin
\textbf{Ōðer} remains strong after the article. It syncopates before vowel-initial \textbf{-es}, but retains \textbf{e} before consonantal \textbf{-re}. Feminine nominative \textbf{ōðer/\allowbreak{}ōðeru} are documented alternatives to selected \textbf{ōðru}. Genitive \textbf{niht} remains endingless; the phrase’s construction and agreement establish its case.
% mdparagraph:P061:end

\end{minipage}\par\addvspace{4pt}

\par\begin{minipage}{\linewidth}\raggedright
% item: NEW-U03-O02
% key: KEY-NEW-U03-O02
% mdparagraph:P062:begin
\textbf{O2.} Complete weak \textbf{ðridda}, “third”:
% mdparagraph:P062:end

% mdtable:T18:begin
\par\begingroup
\fontsize{10.5}{13.2}\selectfont
\setlength{\LTleft}{0pt}\setlength{\LTright}{\fill}
\noindent
\begin{tabular}{@{}>{\raggedright\arraybackslash}p{27.74mm}>{\raggedright\arraybackslash}p{21.25mm}>{\raggedright\arraybackslash}p{21.25mm}>{\raggedright\arraybackslash}p{21.25mm}>{\raggedright\arraybackslash}p{21.25mm}@{}}
\toprule
\textbf{Gender and number} & \textbf{Nominative} & \textbf{Genitive} & \textbf{Dative} & \textbf{Accusative} \\
\midrule
Masculine singular & ðridda & ðriddan & ðriddan & ðriddan \\
Neuter singular & ðridde & ðriddan & ðriddan & ðridde \\
Feminine singular & ðridde & ðriddan & ðriddan & ðriddan \\
Masculine plural & ðriddan & ðridra & ðriddum & ðriddan \\
Neuter plural & ðriddan & ðridra & ðriddum & ðriddan \\
Feminine plural & ðriddan & ðridra & ðriddum & ðriddan \\
\bottomrule
\end{tabular}
\par\endgroup

% mdtable:T18:end
% mdparagraph:P063:begin
This ordinal is weak with or without a preceding article. Selected genitive plural \textbf{ðridra} simplifies \textbf{dd} before another consonant; dative plural \textbf{ðriddum} retains it. Do not copy the strong exception \textbf{ōðer} into the other ordinal paradigms.
% mdparagraph:P063:end

\end{minipage}\par\addvspace{4pt}

\par\begin{minipage}{\linewidth}\raggedright
% item: NEW-U03-O03
% key: KEY-NEW-U03-O03
% mdparagraph:P064:begin
\textbf{O3.} \textbf{Ðrīe fuglas libbað on ðǣm ǣrestan hofe. Se ealda mann sealde ðǣm fīftan fugle męte.} \textbf{Ǣrestan} agrees with neuter dative singular \textbf{hofe}, a location. \textbf{Fīftan} agrees with masculine dative singular \textbf{fugle}, the recipient. Both ordinals are weak. The first clause is present plural; the second is past singular. Here \textbf{on} is locative, not an untaught temporal construction.
% mdparagraph:P064:end

\end{minipage}\par\addvspace{4pt}

\reserve{7\baselineskip}
% mdparagraph:P065:begin
\section*{Demonstratives}
% mdparagraph:P065:end

\par\begin{minipage}{\linewidth}\raggedright
% item: NEW-U03-D01
% key: KEY-NEW-U03-D01
% mdparagraph:P066:begin
\textbf{D1.} Complete selected \textbf{ðes}, “this/\allowbreak{}these”:
% mdparagraph:P066:end

% mdtable:T19:begin
\par\begingroup
\fontsize{10.5}{13.2}\selectfont
\setlength{\LTleft}{0pt}\setlength{\LTright}{\fill}
\noindent
\begin{tabular}{@{}>{\raggedright\arraybackslash}p{19.22mm}>{\raggedright\arraybackslash}p{25.09mm}>{\raggedright\arraybackslash}p{21.44mm}>{\raggedright\arraybackslash}p{23.17mm}>{\raggedright\arraybackslash}p{23.84mm}@{}}
\toprule
\textbf{Case} & \textbf{Masculine singular} & \textbf{Neuter singular} & \textbf{Feminine singular} & \textbf{Plural, all genders} \\
\midrule
Nominative & ðes & ðis & ðēos & ðās \\
Genitive & ðisses & ðisses & ðisse & ðissa \\
Dative & ðissum & ðissum & ðisse & ðissum \\
Accusative & ðisne & ðis & ðās & ðās \\
\bottomrule
\end{tabular}
\par\endgroup

% mdtable:T19:end
% mdparagraph:P067:begin
Keep short \textbf{e} in selected masculine \textbf{ðes} and long \textbf{ē} in feminine \textbf{ðēos}. These are the four requested cases; the distinct instrumental reference forms are not production targets here.
% mdparagraph:P067:end

\end{minipage}\par\addvspace{4pt}

\par\begin{minipage}{\linewidth}\raggedright
% item: NEW-U03-D02
% key: KEY-NEW-U03-D02
% mdparagraph:P068:begin
\textbf{D2.} \textbf{Ðis trēo is micel. Ðæt trēo is lȳtel. Ðās fuglas libbað on ðǣm hofe.} Neuter singular \textbf{ðis/\allowbreak{}ðæt} distinguish this tree from that tree; plural \textbf{ðās} identifies these birds. Dative \textbf{ðǣm} identifies that dwelling after locative \textbf{on}. The context assigns the pointing “that” sense to forms that can also serve as articles. Predicate \textbf{micel/\allowbreak{}lȳtel} are strong.
% mdparagraph:P068:end

\end{minipage}\par\addvspace{4pt}

\par\begin{minipage}{\linewidth}\raggedright
% item: NEW-U03-D03
% key: KEY-NEW-U03-D03
% mdparagraph:P069:begin
\textbf{D3.} \textbf{Ic sęlle ðissum męnn sǣd. Ic hīere ðās fuglas.} \textbf{Ðissum męnn} is masculine dative singular, the recipient. \textbf{Ðās fuglas} is masculine accusative plural, the object heard. The endings agree within each noun phrase; nominative \textbf{ðes mann} would not satisfy the giving construction.
% mdparagraph:P069:end

\end{minipage}\par\addvspace{4pt}

\par\begin{minipage}{\linewidth}\raggedright
% item: NEW-U03-D04
% key: KEY-NEW-U03-D04
% mdparagraph:P070:begin
\textbf{D4.} \textbf{Ðes gōda mann; ðās miclan trēowu; gōde męnn.} The first two adjectives are weak after demonstratives: masculine nominative singular \textbf{gōda} and neuter nominative plural \textbf{miclan}. Bare indefinite \textbf{gōde męnn}, “good men”, has a strong masculine nominative plural adjective. No demonstrative is needed in the last phrase.
% mdparagraph:P070:end

\end{minipage}\par\addvspace{4pt}

\reserve{7\baselineskip}
% mdparagraph:P071:begin
\section*{Controlled reading: Two dwellings and three birds}
% mdparagraph:P071:end

\par\begin{minipage}{\linewidth}\raggedright
% item: NEW-U03-L01
% key: KEY-NEW-U03-L01
% mdparagraph:P072:begin
\textbf{L1.} A translation from the complete Old English account:
% mdparagraph:P072:end

% mdparagraph:P073:begin
“The good men built two dwellings. One dwelling was large, and one dwelling was small. The large dwelling was larger than the small dwelling. They gave the large dwelling to the old man. He had three birds. He gave his birds seed and water. The birds lived in the large dwelling. The man said: ‘My birds are good. Our work is good.’”
% mdparagraph:P073:end

% mdparagraph:P074:begin
“Made” for \textbf{geworhton}, “great/\allowbreak{}little” for \textbf{micel/\allowbreak{}lȳtel}, and “houses/\allowbreak{}homes” for \textbf{hofu} are defensible in this setting. The builders give the dwelling; the old man receives it and then gives seed and water to his birds. By the stated orientation he is one of the builders: \textbf{mīne} refers to his birds, while \textbf{ūre} includes the group’s work. Do not change the two quoted present statements to past merely because \textbf{sægde} is past.
% mdparagraph:P074:end

\end{minipage}\par\addvspace{4pt}

\par\begin{minipage}{\linewidth}\raggedright
% item: NEW-U03-L02
% key: KEY-NEW-U03-L02
% mdparagraph:P075:begin
\textbf{L2.} The counted things and analyses are:
% mdparagraph:P075:end

% mdtable:T20:begin
\par\begingroup
\fontsize{10.5}{13.2}\selectfont
\setlength{\LTleft}{0pt}\setlength{\LTright}{\fill}
\noindent
\begin{tabular}{@{}>{\raggedright\arraybackslash}p{24.79mm}>{\raggedright\arraybackslash}p{37.35mm}>{\raggedright\arraybackslash}p{56.24mm}@{}}
\toprule
\textbf{Occurrence} & \textbf{Counted noun} & \textbf{Numeral analysis} \\
\midrule
twā hofu & two dwellings & neuter accusative plural: object of geworhton \\
first ān hof & one dwelling, the large one & neuter nominative singular: subject of wæs \\
second ān hof & one dwelling, the small one & neuter nominative singular: subject of wæs \\
ðrīe fuglas & three birds & masculine accusative plural: object of hæfde \\
\bottomrule
\end{tabular}
\par\endgroup

% mdtable:T20:end
% mdparagraph:P076:begin
The two identical \textbf{ān} forms have the same grammatical analysis but identify different dwellings. \textbf{Twā/\allowbreak{}ðrīe} agree with plural nouns; their spelling alone does not distinguish nominative from accusative.
% mdparagraph:P076:end

\end{minipage}\par\addvspace{4pt}

\par\begin{minipage}{\linewidth}\raggedright
% item: NEW-U03-L03
% key: KEY-NEW-U03-L03
% mdparagraph:P077:begin
\textbf{L3.} In \textbf{Mīne fuglas}, the owner is the speaking old man, first-person singular; the possessed noun is masculine nominative plural \textbf{fuglas}. \textbf{Mīne} has the corresponding agreeing \textbf{-e} form. In \textbf{Ūre weorc}, the owners are all the builders including the speaker, first-person plural; the possessed noun is neuter nominative singular \textbf{weorc}. \textbf{Ūre} is the selected singular neuter form of that possessive. Final \textbf{-e} therefore does not identify one unique gender or number, and it does not itself count the owners.
% mdparagraph:P077:end

\end{minipage}\par\addvspace{4pt}

\par\begin{minipage}{\linewidth}\raggedright
% item: NEW-U03-L04
% key: KEY-NEW-U03-L04
% mdparagraph:P078:begin
\textbf{L4.} The actual adjective occurrences are:
% mdparagraph:P078:end

% mdtable:T21:begin
\par\begingroup
\fontsize{10.5}{13.2}\selectfont
\setlength{\LTleft}{0pt}\setlength{\LTright}{\fill}
\noindent
\begin{tabular}{@{}>{\raggedright\arraybackslash}p{34.47mm}>{\raggedright\arraybackslash}p{33.89mm}>{\raggedright\arraybackslash}p{50.01mm}@{}}
\toprule
\textbf{Kind} & \textbf{Form and location} & \textbf{Agreement} \\
\midrule
Strong predicate & micel; lȳtel, second sentence & each with its nominative neuter singular hof \\
Strong predicate & gōde, first quoted sentence & with nominative masculine plural fuglas \\
Strong predicate & gōd, second quoted sentence & with nominative neuter singular weorc \\
Weak attributive & gōdan, first sentence & with nominative masculine plural męnn \\
Weak attributive & micle; lȳtle, third sentence & each with nominative neuter singular hof \\
Weak attributive & micle, fourth sentence & with accusative neuter singular hof \\
Weak attributive & ealdan, fourth sentence & with dative masculine singular męnn \\
Weak attributive & miclan, seventh sentence & with dative neuter singular hofe \\
Weak comparative predicate & māre, third sentence & with the first nominative neuter singular hof \\
\bottomrule
\end{tabular}
\par\endgroup

% mdtable:T21:end
% mdparagraph:P079:begin
The comparative remains weak in predication. The same spelling \textbf{micle} serves different cases in the third and fourth sentences; clause roles identify the cells. \textbf{Mīne/\allowbreak{}ūre} are agreeing possessives, not additional weak attributive or predicate adjectives in this requested list.
% mdparagraph:P079:end

\end{minipage}\par\addvspace{4pt}

\par\begin{minipage}{\linewidth}\raggedright
% item: NEW-U03-L05
% key: KEY-NEW-U03-L05
% mdparagraph:P080:begin
\textbf{L5.} \textbf{Ðā gōdan męnn gewyrceað twā hofu. Ān hof is micel, and ān hof is lȳtel. Ðæt micle hof is māre ðonne ðæt lȳtle hof. Hīe sęllað ðæt micle hof ðǣm ealdan męnn.}
% mdparagraph:P080:end

% mdparagraph:P081:begin
The changes are \textbf{geworhton → gewyrceað}, each \textbf{wæs → is}, and \textbf{sealdon → sęllað}. Noun-phrase cases, adjective endings, degree and referents remain unchanged. The plural builders still govern the plural active verbs; the singular dwellings govern \textbf{is}.
% mdparagraph:P081:end

\end{minipage}\par\addvspace{4pt}

\reserve{7\baselineskip}
% mdparagraph:P082:begin
\section*{Source reading: the fourth day}
% mdparagraph:P082:end

\par\begin{minipage}{\linewidth}\raggedright
% item: NEW-U03-R01
% key: KEY-NEW-U03-R01
% mdparagraph:P083:begin
\textbf{R1.} A close translation from the complete supplied source:
% mdparagraph:P083:end

% mdparagraph:P084:begin
“God then said, ‘Let there now be lights in the firmament of heaven, and let them separate day and night, and let them be for signs, for seasons, for days, and for years. And let them shine in the firmament of heaven and illuminate the earth’; it then came to pass thus. And God made two great lights: the greater light for the illumination of the day, and the lesser light for the illumination of the night; and he made stars. And he set them in heaven, so that they might shine over the earth, rule over the day and the night, and separate light and darkness; God then saw that it was good. And evening and morning came to pass, the fourth day.” [3]
% mdparagraph:P084:end

% mdparagraph:P085:begin
The opening \textbf{sōðlice} is a narrative connective here, glossed by Cook’s Latin \textbf{autem}; the English opening absorbs it. “Let them serve as signs…” is a reasonable idiomatic paraphrase of \textbf{bēon tō tācnum}. “It was then so” can render the fulfillment formula, and “there came evening and morning, the fourth day” the closing formula. Preserve the wish, fulfillment, and all three later purposes: shining, ruling, and separating.
% mdparagraph:P085:end

% mdparagraph:P086:begin
The two lights remain unnamed in this source. \textbf{Heofonan} is genitive in “the firmament of heaven”, dative in “in heaven”; the supplied omitted subject of \textbf{gesętte} is God/\allowbreak{}he. \textbf{Gīemden} governs the genitives for day and night, whereas \textbf{tō} governs the illumination phrases’ dative heads. The source’s \textbf{męrgen} is the supplied variant of \textbf{morgen}. These supports may be consulted; neither a unique case for \textbf{ofer eorðan} nor unaided mastery of the supplied moods and resultative constructions is required.
% mdparagraph:P086:end

\end{minipage}\par\addvspace{4pt}

\par\begin{minipage}{\linewidth}\raggedright
% item: NEW-U03-R02
% key: KEY-NEW-U03-R02
% mdparagraph:P087:begin
\textbf{R2.} All three words in \textbf{twā miclu lēoht} are \textbf{neuter accusative plural}, the direct object of \textbf{geworhte}:
% mdparagraph:P087:end

% mdtable:T22:begin
\par\begingroup
\fontsize{10.5}{13.2}\selectfont
\setlength{\LTleft}{0pt}\setlength{\LTright}{\fill}
\noindent
\begin{tabular}{@{}>{\raggedright\arraybackslash}p{22.74mm}>{\raggedright\arraybackslash}p{22.74mm}>{\raggedright\arraybackslash}p{72.89mm}@{}}
\toprule
\textbf{Source form} & \textbf{Headword} & \textbf{Analysis} \\
\midrule
twā & twēgen & cardinal two, agreeing with plural neuter lēoht \\
miclu & micel & strong adjective accusative neuter plural \\
lēoht & lēoht & heavy neuter count noun, accusative plural with zero ending \\
\bottomrule
\end{tabular}
\par\endgroup

% mdtable:T22:end
% mdparagraph:P088:begin
\textbf{Micel} has selected syncopated \textbf{micl-} before \textbf{-u}. The adjective is strong in this numeral phrase, which has no preceding article/\allowbreak{}demonstrative. The noun means two lights or luminaries, not an uncounted quantity of light. Preserve the source’s exact \textbf{miclu}.
% mdparagraph:P088:end

\end{minipage}\par\addvspace{4pt}

\par\begin{minipage}{\linewidth}\raggedright
% item: NEW-U03-R03
% key: KEY-NEW-U03-R03
% mdparagraph:P089:begin
\textbf{R3.} \textbf{Māre} is the comparative of \textbf{micel}, “great/\allowbreak{}large”; \textbf{lǣsse} of \textbf{lȳtel}, “little/\allowbreak{}small”. Both are weak \textbf{accusative neuter singular}, agreeing with the respective singular \textbf{lēoht} in apposition to the two objects made. Their neuter nominative spellings are identical, but the larger object construction fixes accusative here.
% mdparagraph:P089:end

% mdparagraph:P090:begin
The greater light is for the illumination of the day; the lesser is for the illumination of the night. Those purposes follow the supplied \textbf{tō ... līehtinge} construction. The passage does not require naming either light.
% mdparagraph:P090:end

\end{minipage}\par\addvspace{4pt}

\par\begin{minipage}{\linewidth}\raggedright
% item: NEW-U03-R04
% key: KEY-NEW-U03-R04
% mdparagraph:P091:begin
\textbf{R4.} The noun phrases are \textbf{ðæt māre lēoht} and \textbf{ðæt lǣsse lēoht}. Each \textbf{ðæt} is the neuter accusative singular article, agreeing with its light.
% mdparagraph:P091:end

% mdparagraph:P092:begin
In \textbf{ðæt hīe scinen}, supplied \textbf{ðæt} introduces purpose: “so that they might shine”. In \textbf{ðæt hit gōd wæs}, it introduces the content of what God saw: “that it was good”. Neither conjunction agrees with a noun or supplies an article case. Their identical spelling does not make these four occurrences one construction.
% mdparagraph:P092:end

\end{minipage}\par\addvspace{4pt}

\par\begin{minipage}{\linewidth}\raggedright
% item: NEW-U03-R05
% key: KEY-NEW-U03-R05
% mdparagraph:P093:begin
\textbf{R5.} Both \textbf{geworhte} forms are \textbf{third-person singular preterite indicative} of \textbf{gewyrcean}. The first has subject \textbf{God} and object \textbf{twā miclu lēoht}, followed by the two explanatory singular light phrases. The second has subject \textbf{hē} and object \textbf{steorran}, “stars”, which precedes the subject and verb.
% mdparagraph:P093:end

% mdparagraph:P094:begin
The supplied \textbf{Bēon} and long-vowel \textbf{scīnen} express plural wishes: “let there be” and “let them shine”. They are present optative/\allowbreak{}subjunctive forms of different verbs, not members of the taught weak indicative paradigm of \textbf{gewyrcean}. The local help supplies this distinction; independent production of those moods is not tested. Later short-vowel \textbf{scinen} is a separately supplied past optative in the purpose clause.
% mdparagraph:P094:end

\end{minipage}\par\addvspace{4pt}

\par\begin{minipage}{\linewidth}\raggedright
% item: NEW-U03-R06
% key: KEY-NEW-U03-R06
% mdparagraph:P095:begin
\textbf{R6.} \textbf{Gōd} is strong positive \textbf{nominative neuter singular}, agreeing predicatively with \textbf{hit} in the supplied content clause. A predicate adjective does not become weak because the clause begins with conjunction \textbf{ðæt}.
% mdparagraph:P095:end

% mdparagraph:P096:begin
\textbf{Fēorða} is weak ordinal \textbf{nominative masculine singular}, agreeing with \textbf{dæg} in \textbf{se fēorða dæg}, “the fourth day”. The selected ordinals other than \textbf{ōðer} are weak. The larger closing resultative formula remains supplied; its interpretation does not require inventing a new independent agreement rule.
% mdparagraph:P096:end

\end{minipage}\par\addvspace{4pt}

\reserve{7\baselineskip}
% mdparagraph:P097:begin
\section*{Review}
% mdparagraph:P097:end

\par\begin{minipage}{\linewidth}\raggedright
% item: NEW-U03-X01
% key: KEY-NEW-U03-X01
% mdparagraph:P098:begin
\textbf{X1.} \textbf{Ðes gōda mann sęlð mīnum fuglum sǣd. Mīne fuglas libbað on his lȳtlum hofe. Ðæt hof is māre ðonne ðæt trēo.}
% mdparagraph:P098:end

% mdparagraph:P099:begin
The birds are dative recipients in the first clause, nominative subjects in the second. \textbf{Mīnum/\allowbreak{}mīne} change accordingly; \textbf{his} remains invariant. \textbf{Lȳtlum} is the selected strong dative neuter singular adjective after the possessor. Weak \textbf{lȳtlan} is also a defensible adjective choice there: neither must be rejected merely because the other is selected. Both compared nouns are nominative, and \textbf{māre} is weak neuter singular. Source alternative \textbf{sęleð} is valid for selected \textbf{sęlð}.
% mdparagraph:P099:end

\end{minipage}\par\addvspace{4pt}

\par\begin{minipage}{\linewidth}\raggedright
% item: NEW-U03-X02
% key: KEY-NEW-U03-X02
% mdparagraph:P100:begin
\textbf{X2.} \textbf{Se ōðer dæg wæs gōd. Mīne fuglas libbað on ðǣm miclan hofe. Se fugol is lǣssa.}
% mdparagraph:P100:end

% mdparagraph:P101:begin
The four repairs are \textbf{ōðra → ōðer}, because this ordinal remains strong; \textbf{Mīn → Mīne}, for masculine nominative plural agreement; \textbf{micel → miclan}, for weak dative neuter singular after the article; and \textbf{lǣss → lǣssa}, for a weak masculine nominative singular comparative. \textbf{Gōd} is already the correct strong predicate, so it does not change.
% mdparagraph:P101:end

\end{minipage}\par\addvspace{4pt}

\par\begin{minipage}{\linewidth}\raggedright
% item: NEW-U03-X03
% key: KEY-NEW-U03-X03
% mdparagraph:P102:begin
\textbf{X3.} One four-sentence model:
% mdparagraph:P102:end

% mdparagraph:P103:begin
\textbf{Ān hof is micel, and ān hof is lȳtel. Ðæt micle hof is māre ðonne ðæt lȳtle hof. Eald mann hæfð ðrīe fuglas. Hē sægð: “Ic sęlle mīnum fuglum męte.”}
% mdparagraph:P103:end

% mdparagraph:P104:begin
“One dwelling is large, and one dwelling is small. The large dwelling is larger than the small dwelling. An old man has three birds. He says: ‘I give my birds food.’”
% mdparagraph:P104:end

% mdparagraph:P105:begin
Other connected solutions are acceptable within the taught vocabulary and constructions. The essentials are two distinguishable dwellings of unequal size, a comparison, an old man with three birds, and a quotation in which he gives those birds food. In the model \textbf{mīnum} is an agreeing possessive with dative plural recipients \textbf{fuglum}; \textbf{męte} is the accusative thing given. The first-person \textbf{ic/\allowbreak{}mīn-} belongs to the speaking man. Merely reporting that he gives food, without the requested quotation and agreeing possessive, would miss those targets.
% mdparagraph:P105:end

\end{minipage}\par\addvspace{4pt}

\reserve{7\baselineskip}
% mdparagraph:P106:begin
\section*{Cook I.4, items 1–17}
% mdparagraph:P106:end

% mdparagraph:P107:begin
The original item numbers are retained. Noun-phrase heads use the requested nominative convention; dependent genitives retain their own cases. A contextual accusative phrase can be valid Old English while missing that stated target. Put dependent genitives before their heads as requested. Equivalent postposed genitives are valid Old English, but miss this exercise’s positioning convention.
% mdparagraph:P107:end

\par\begin{minipage}{\linewidth}\raggedright
% item: EX-I-04-01
% key: KEY-EX-I-04-01
% mdparagraph:P108:begin
\textbf{1.} \textbf{Ðā miclan hwalas.} “The great whales.” The adjective is weak nominative masculine plural after \textbf{ðā}; \textbf{hwalas} has the taught \textbf{æ/\allowbreak{}a} noun-stem alternation.
% mdparagraph:P108:end

\end{minipage}\par\addvspace{4pt}

\par\begin{minipage}{\linewidth}\raggedright
% item: EX-I-04-02
% key: KEY-EX-I-04-02
% mdparagraph:P109:begin
\textbf{2.} \textbf{Sēo eorðe wæs īdel and ǣmtig.} Both adjectives are strong nominative feminine singular predicates. \textbf{Ǣmtig and īdel} is also defensible: Cook’s vocabulary gives both for empty/\allowbreak{}void. The subject’s article does not make these predicate adjectives weak.
% mdparagraph:P109:end

\end{minipage}\par\addvspace{4pt}

\par\begin{minipage}{\linewidth}\raggedright
% item: EX-I-04-03
% key: KEY-EX-I-04-03
% mdparagraph:P110:begin
\textbf{3.} \textbf{Ðā ðing wǣron gōd.} “The things were good.” \textbf{Ðing} is nominative neuter plural with zero ending; plural \textbf{wǣron} agrees with it. Strong neuter plural predicate \textbf{gōd} likewise has no added \textbf{-u}. Masculine \textbf{gōde} would not match the noun’s gender.
% mdparagraph:P110:end

\end{minipage}\par\addvspace{4pt}

\par\begin{minipage}{\linewidth}\raggedright
% item: EX-I-04-04
% key: KEY-EX-I-04-04
% mdparagraph:P111:begin
\textbf{4.} \textbf{Ðā æppelbǣran trēowu.} The adjective is weak nominative neuter plural; \textbf{trēowu} is the taught plural of \textbf{trēo}. The attested compound \textbf{æppelbǣre} supplies the meaning “fruit-bearing”.
% mdparagraph:P111:end

\end{minipage}\par\addvspace{4pt}

\par\begin{minipage}{\linewidth}\raggedright
% item: EX-I-04-05
% key: KEY-EX-I-04-05
% mdparagraph:P112:begin
\textbf{5.} \textbf{Ðā miclan lēoht.} “The great lights.” The weak adjective and article accompany a nominative neuter plural count noun. This is the new luminary sense, not the earlier uncounted “light”. Zero noun ending does not make the phrase singular.
% mdparagraph:P112:end

\end{minipage}\par\addvspace{4pt}

\par\begin{minipage}{\linewidth}\raggedright
% item: EX-I-04-06
% key: KEY-EX-I-04-06
% mdparagraph:P113:begin
\textbf{6.} \textbf{Se lȳtla steorra.} The adjective is weak nominative masculine singular, with selected stem \textbf{lȳtl-}. \textbf{Steorra} is masculine despite its final \textbf{-a}; neuter/\allowbreak{}feminine \textbf{lȳtle} would not be the selected agreement.
% mdparagraph:P113:end

\end{minipage}\par\addvspace{4pt}

\par\begin{minipage}{\linewidth}\raggedright
% item: EX-I-04-07
% key: KEY-EX-I-04-07
% mdparagraph:P114:begin
\textbf{7.} \textbf{Se lǣssa fugol.} “The smaller bird.” Comparative \textbf{lǣssa} is weak nominative masculine singular; the comparative is weak even without an article.
% mdparagraph:P114:end

\end{minipage}\par\addvspace{4pt}

\par\begin{minipage}{\linewidth}\raggedright
% item: EX-I-04-08
% key: KEY-EX-I-04-08
% mdparagraph:P115:begin
\textbf{8.} \textbf{Ðæt māre trēo.} “The larger tree.” Comparative \textbf{māre} is weak nominative neuter singular. Masculine \textbf{māra} would not agree with \textbf{trēo}.
% mdparagraph:P115:end

\end{minipage}\par\addvspace{4pt}

\par\begin{minipage}{\linewidth}\raggedright
% item: EX-I-04-09
% key: KEY-EX-I-04-09
% mdparagraph:P116:begin
\textbf{9.} \textbf{Ðæs gōdan mannes līf.} \textbf{Ðæs gōdan mannes} is genitive masculine singular, “of the good man”; \textbf{līf} is the nominative neuter head. \textbf{Gōdan} agrees with the man, not with life. \textbf{Ðæt līf ðæs gōdan mannes} is valid Old English, but its postposed genitive misses the requested positioning convention.
% mdparagraph:P116:end

\end{minipage}\par\addvspace{4pt}

\par\begin{minipage}{\linewidth}\raggedright
% item: EX-I-04-10
% key: KEY-EX-I-04-10
% mdparagraph:P117:begin
\textbf{10.} \textbf{Ðǣre seofoðan niht anginn.} \textbf{Seofoðan} is weak genitive feminine singular with \textbf{niht}; \textbf{anginn} is the nominative neuter head, “beginning”. The ordinal is weak in its own right. Endingless \textbf{niht} is genitive here.
% mdparagraph:P117:end

\end{minipage}\par\addvspace{4pt}

\par\begin{minipage}{\linewidth}\raggedright
% item: EX-I-04-11
% key: KEY-EX-I-04-11
% mdparagraph:P118:begin
\textbf{11.} \textbf{Ðæs miclan sǣs wæteru.} “The waters of the great sea.” \textbf{Ðæs miclan sǣs} is genitive masculine singular; \textbf{wæteru} is nominative neuter plural. The adjective describes the sea, not the waters. Preserve the taught singular genitive \textbf{sǣs}.
% mdparagraph:P118:end

\end{minipage}\par\addvspace{4pt}

\par\begin{minipage}{\linewidth}\raggedright
% item: EX-I-04-12
% key: KEY-EX-I-04-12
% mdparagraph:P119:begin
\textbf{12.} \textbf{Se ōðer dæg and sēo ōðru niht.} Both heads are nominative singular. \textbf{Ōðer} stays strong after the article; selected feminine \textbf{ōðru} agrees with \textbf{niht}. \textbf{Sēo ōðer niht} and \textbf{sēo ōðeru niht} are documented alternatives. The original \textbf{(80)} is Cook’s grammatical reference, not part of the phrase to translate.
% mdparagraph:P119:end

\end{minipage}\par\addvspace{4pt}

\par\begin{minipage}{\linewidth}\raggedright
% item: EX-I-04-13
% key: KEY-EX-I-04-13
% mdparagraph:P120:begin
\textbf{13.} \textbf{Ðæs ōðres gēares ðridda dæg.} The limiting genitive \textbf{ðæs ōðres gēares} is neuter singular, “of the second year”; \textbf{ðridda dæg} is nominative masculine singular, “the third day”. \textbf{Ōðres} is strong, while \textbf{ðridda} is weak even without an article immediately before it. \textbf{Se ðridda dæg ðæs ōðres gēares} is valid Old English, but its postposed genitive misses the requested positioning convention.
% mdparagraph:P120:end

\end{minipage}\par\addvspace{4pt}

\par\begin{minipage}{\linewidth}\raggedright
% item: EX-I-04-14
% key: KEY-EX-I-04-14
% mdparagraph:P121:begin
\textbf{14.} \textbf{Ðæs fīftan mannes męte.} The ordinal is weak genitive masculine singular with \textbf{mannes}. \textbf{Męte}, “food”, is the nominative masculine head; its final \textbf{-e} does not make it dative here.
% mdparagraph:P121:end

\end{minipage}\par\addvspace{4pt}

\par\begin{minipage}{\linewidth}\raggedright
% item: EX-I-04-15
% key: KEY-EX-I-04-15
% mdparagraph:P122:begin
\textbf{15.} \textbf{Ðæs siextan ǣfenes lēoht.} \textbf{Siextan} agrees with genitive neuter singular \textbf{ǣfenes}. The head \textbf{lēoht} is nominative neuter singular in the mass sense “light”, not a plurality of luminaries.
% mdparagraph:P122:end

\end{minipage}\par\addvspace{4pt}

\par\begin{minipage}{\linewidth}\raggedright
% item: EX-I-04-16
% key: KEY-EX-I-04-16
% mdparagraph:P123:begin
\textbf{16.} \textbf{Ðāra styriendlicra gesceafta weorc.} The phrase \textbf{ðāra styriendlicra gesceafta} is genitive feminine plural; \textbf{weorc} is the nominative neuter plural head. Selected weak-adjective genitive \textbf{styriendlicra} has \textbf{-ra}. \textbf{Styriendlic}, “moving”, is the taught lexical adjective; no independently recalled participle is required. Documented ordinary weak-genitive alternatives remain valid, but the model shows the selected ending.
% mdparagraph:P123:end

\end{minipage}\par\addvspace{4pt}

\par\begin{minipage}{\linewidth}\raggedright
% item: EX-I-04-17
% key: KEY-EX-I-04-17
% mdparagraph:P124:begin
\textbf{17.} \textbf{Ðæs ǣrestan dæges morgen.} \textbf{Ǣrestan} is weak genitive masculine singular with \textbf{dæges}; \textbf{morgen} is the nominative masculine head. The selected ordinal \textbf{ǣresta} means “first”. The source reading’s supplied \textbf{męrgen} is a documented variant, but \textbf{morgen} remains the ordinary selected production form.
% mdparagraph:P124:end

\end{minipage}\par\addvspace{4pt}

\clearpage
% mdparagraph:P125:begin
\section*{Source notes}
% mdparagraph:P125:end

\begingroup\fontsize{10}{12.8}\selectfont
% mdparagraph:P126:begin
[1] Cook, \textbf{A First Book in Old English}, second edition (1895), §§55–80, pp. 38–48, and §§83–85, pp. 49–51, provides the adjective, degree, numeral, possessive and demonstrative systems; §182, p. 100, treats adverbial modification. One-based PDF pages are printed pages +16. Cook §60 admits ordinary weak genitive-plural alternatives. Wright and Wright, \textbf{Old English Grammar}, second edition (1914), §§223, 259, 421–435, 440, 443–445, 448–452, 464 and 555–556 corroborates the checked stem changes, comparisons, possessives, numerals and adverbs. Selected comparative genitive plurals are applications of those stated patterns, not claims that Cook prints a full table for each word.
% mdparagraph:P126:end

% mdparagraph:P127:begin
[2] R. D. Fulk, \textbf{An Introductory Grammar of Old English} (2014), corrected open-access reprint June 2020, §35, p. 22/\allowbreak{}PDF 34, prints a comparison with nominative \textbf{iċ} after \textbf{þonne}, attributed to \textbf{Beowulf} 469. The present exercises adapt that bounded comparison between subjects; they do not teach a universal nominative rule after every \textbf{ðonne}.
% mdparagraph:P127:end

% mdparagraph:P128:begin
[3] Cook, Reader I, \textbf{The Creation of the World}, pp. 125–126/\allowbreak{}PDF 141–142: p. 125 lines 8–15 and p. 126 lines 1–4. Only marginal 10 and 15 are printed within the selected span; the remaining line positions are editorial counts. The translation follows the complete checked Old English and its separate apparatus, including p. 125 note 15 (\textbf{autem}), note 17 (genitive \textbf{heofonan}), note 23 (present optative), note 29 (omitted subject), and p. 126 notes 1–3 (past optative, governed genitive, ordinal). The source help also draws on Cook §§53, 138, 156.f, 166, 189, 193.a, 194 and 196.f. These notes support the local reading; they confer no additional production scope.
% mdparagraph:P128:end

\endgroup
\end{document}
